%=============================================================================
% Wave 3, Iteration 7: Materials Science Update
% AGENT: MATERIALS SCIENCE (W3i7)
%
% W3i6 KEY FINDINGS CARRIED FORWARD:
%   - Passive FOM: Nb dominates (inverts active ranking)
%   - Si3N4 optimal MEMS material (dissipation dilution -> Q_mech > 10^9)
%   - Three-track roadmap: Passive/Nb, MEMS/Si3N4, Exotic/YBCO monitoring
%   - P11 Phase I: SQMS 4-cavity at Fermilab, $2-4M, 2027-2028
%   - Active FOM unchanged under two-component decomposition
%   - YBCO demoted for passive detection (d-wave nodes)
%
% W3i7 MANDATE:
%   1. Si3N4 MEMS: Complete fabrication spec (dimensions, stress,
%      resonant frequency, Q_mech). Cross-reference state-of-art.
%      Is 10^9 Q_mech at nanogram mass demonstrated or projected?
%   2. Nb cavity optimization: For SQMS 4-cavity Phase I, what Nb
%      preparation maximizes passive sensitivity? Specific Q0 measurement
%      that constitutes indirect evidence of psi-coupling?
%   3. Novel materials: Superconducting nanowires, topological surface
%      states, superfluid He-3, dilute Bose gas -- any exceptional
%      passive FOM?
%   4. Cryogenic engineering: Phase I-III cost and power budget.
%      Does this change the economics?
%
% Date: March 2026
%=============================================================================

\documentclass[11pt]{article}
\usepackage{amsmath,amssymb,amsthm}
\usepackage{geometry}
\geometry{margin=1in}
\usepackage{booktabs}
\usepackage{longtable}
\usepackage{array}
\usepackage{enumitem}
\usepackage{hyperref}
\usepackage{xcolor}
\usepackage{tcolorbox}
\usepackage{float}
\usepackage{siunitx}

\newtheorem{theorem}{Theorem}[section]
\newtheorem{proposition}[theorem]{Proposition}
\newtheorem{corollary}[theorem]{Corollary}
\newtheorem{lemma}[theorem]{Lemma}
\theoremstyle{definition}
\newtheorem{definition}[theorem]{Definition}
\newtheorem{finding}[theorem]{Finding}
\newtheorem{correction}[theorem]{Correction}
\theoremstyle{remark}
\newtheorem{remark}[theorem]{Remark}

\newcommand{\ee}[1]{\times 10^{#1}}
\newcommand{\Meff}{M_{\rm eff}}
\newcommand{\Opsi}{\omega_\psi}
\newcommand{\Qpsi}{Q_\psi}
\newcommand{\dpsi}{\delta\psi_{\rm EM}}
\newcommand{\psigrav}{\psi_{\rm grav}}
\newcommand{\psiEM}{\delta\psi_{\rm EM}}
\newcommand{\lambdaone}{|\lambda-1|}
\newcommand{\lambdabound}{1.56\ee{-9}}
\newcommand{\Tc}{T_{\rm c}}
\newcommand{\lambdaL}{\lambda_{\rm L}}
\newcommand{\FOMactive}{\mathrm{FOM}_{\rm active}}
\newcommand{\FOMpassive}{\mathrm{FOM}_{\rm passive}}
\newcommand{\FOMMEMS}{\mathrm{FOM}_{\rm MEMS}}
\newcommand{\lam}{|\lambda-1|}
\newcommand{\Heff}{H_{\rm eff}}
\newcommand{\muG}{\mu_G}
\newcommand{\GN}{G_{\rm N}}
\newcommand{\SNR}{\mathrm{SNR}}
\newcommand{\Qmech}{Q_{\rm mech}}
\newcommand{\confirmed}[1]{\textcolor{blue}{\textbf{CONFIRMED:} #1}}
\newcommand{\corrected}[1]{\textcolor{red}{\textbf{CORRECTED:} #1}}
\newcommand{\caution}[1]{\textcolor{orange}{\textbf{CAUTION:} #1}}
\newcommand{\honest}[1]{\textcolor{violet}{\textbf{HONEST ASSESSMENT:} #1}}
\newcommand{\killed}[1]{\textcolor{red!70!black}{\textbf{KILLED:} #1}}
\newcommand{\verdict}[1]{\textcolor{green!50!black}{\textbf{VERDICT:} #1}}

\title{Wave 3 Iteration 7: Materials Science Update\\
\large Si$_3$N$_4$ MEMS Fabrication Specification,\\
Nb Cavity Optimisation for Phase~I,\\
Novel Material Candidates, and Cryogenic Engineering Budget}
\author{DFD Research Programme --- Materials Science Agent, W3i7}
\date{March 2026}

\begin{document}
\maketitle
\tableofcontents
\newpage

%=============================================================================
\section{Executive Summary}
\label{sec:executive}
%=============================================================================

Wave 3 Iteration 6 established the three-track materials roadmap
(Passive/Nb, MEMS/Si$_3$N$_4$, Exotic/YBCO) and the Phase~I--III
detection programme centred on the SQMS 4-cavity array at Fermilab.
This iteration delivers four deep-dive analyses requested by the
programme:

\begin{tcolorbox}[colback=blue!5!white, colframe=blue!50!black,
  title=\textbf{W3i7 Materials Science Mandate}]
\begin{enumerate}[nosep]
  \item \textbf{Si$_3$N$_4$ MEMS fabrication spec}: Complete dimensions,
    stress, resonant frequency, $\Qmech$. Is $\Qmech = 10^9$ at nanogram
    mass demonstrated or projected?
  \item \textbf{Nb cavity optimisation}: What preparation protocol
    maximises Phase~I passive sensitivity? What $Q_0$ measurement would
    constitute indirect evidence of $\psi$-field coupling?
  \item \textbf{Novel materials}: Superconducting nanowires, topological
    surface states, superfluid He-3, dilute Bose gas---any exceptional
    passive FOM not yet considered?
  \item \textbf{Cryogenic engineering}: Phase~I--III cost and power budget.
    Does cryogenic infrastructure change the programme economics?
\end{enumerate}
\end{tcolorbox}

\noindent\textbf{Key results}:
\begin{itemize}[nosep]
  \item We provide a \textbf{complete Si$_3$N$_4$ MEMS fabrication
    specification} for the Phase~II P2+P11 hybrid: soft-clamped
    phononic crystal membrane, 1~mm $\times$ 1~mm $\times$ 20~nm,
    $\sigma = 1.2$~GPa, $f_0 \approx 1.4$~MHz, $m_{\rm eff} = 0.6$~ng.
    $\Qmech \geq 10^9$ at room temperature is \textbf{demonstrated};
    at millikelvin temperatures $\Qmech > 10^{10}$ is \textbf{projected
    with high confidence} but not yet demonstrated at the required
    mass scale.
  \item The \textbf{optimal Nb preparation for Phase~I} is:
    medium-$T$ N-doping (800$^\circ$C, 20~min, 2/0/2 or 2/6 recipe),
    followed by 5~$\mu$m EP, then low-temperature baking at 120$^\circ$C
    for 48~h, in a sub-nT magnetically shielded environment. We identify
    a specific \textbf{$Q_0$ anomaly signature}---a
    temperature-dependent $\delta Q/Q$ with the spectral shape of the
    DFD $\psi$-mode dispersion---that would constitute circumstantial
    evidence of $\psi$-field coupling.
  \item Of the four novel materials considered, \textbf{superfluid $^3$He}
    emerges as a genuinely new candidate with potential passive FOM
    exceeding Nb for specific geometries (acoustic Q-factors $>10^{10}$
    in torsional oscillators). Superconducting nanowires, topological
    surface states, and dilute Bose gases are assessed and found
    \textbf{not competitive} for the passive programme.
  \item The \textbf{cryogenic power budget} for Phase~I (4-cavity at
    2~K) is approximately 30~kW wall-plug; for Phase~I at 20~mK
    (dilution refrigerator), approximately 50~kW. Phase~III
    (256-cavity) requires 2--5~MW of cryogenic power, comparable to
    a medium-scale accelerator facility. This does \textbf{not} change
    the Phase~I/II economics but makes Phase~III a facility-class
    investment.
\end{itemize}


%=============================================================================
\section{Si$_3$N$_4$ MEMS: Complete Fabrication Specification}
\label{sec:si3n4_spec}
%=============================================================================

\subsection{Context: the P2+P11 hybrid detector requirement}

The P2+P11 hybrid (W3i5 P10/P11, W3i6 MatSci Sec.~6) requires a
mechanical resonator with:
\begin{itemize}[nosep]
  \item $\Qmech \geq 10^9$ to achieve $\lam < 10^{-14}$
    (Section~\ref{sec:sensitivity_recalc});
  \item Effective mass $m_{\rm eff} \sim 0.1$--10~ng (nanogram regime);
  \item Resonant frequency $f_0$ compatible with SQUID readout bandwidth
    ($f_0 \lesssim 10$~MHz);
  \item Cryogenic operation at $T \leq 100$~mK;
  \item SQUID-compatible displacement transduction.
\end{itemize}

\subsection{Design: soft-clamped phononic crystal membrane}
\label{sec:design}

The highest $\Qmech$ demonstrated in Si$_3$N$_4$ devices uses
\textbf{soft-clamped phononic crystal (PnC) membranes}, pioneered by
the groups of Schliesser (Copenhagen/Delft), Aspelmeyer (Vienna), and
Harris (Yale). The key physics:

\begin{enumerate}
  \item \textbf{Dissipation dilution}: In a pre-stressed membrane,
    elastic energy is stored predominantly in the tensile stress field
    rather than in bending deformations at the clamping points. The
    effective quality factor is enhanced by a dilution factor:
    \begin{equation}
      D_{\rm dil} = \frac{\sigma}{\sigma_{\rm crit}}
      = \frac{\sigma L^2}{E h^2}\cdot C_{\rm geom},
      \label{eq:dilution}
    \end{equation}
    where $\sigma$ is the tensile stress, $E$ is Young's modulus, $h$
    is the membrane thickness, $L$ is the characteristic length, and
    $C_{\rm geom}$ is a geometry-dependent factor of order unity.

  \item \textbf{Soft clamping}: A phononic crystal pattern etched into
    the membrane creates a bandgap that prevents vibrational energy from
    reaching the physical clamping points. This eliminates radiation loss
    at the clamps, which otherwise limits $\Qmech$ in conventional
    membranes. The effective dissipation dilution is enhanced by a
    further factor:
    \begin{equation}
      D_{\rm soft} = \frac{n\pi^2}{12}\left(\frac{L}{h}\right)^2
      \cdot \frac{\sigma}{E},
      \label{eq:soft_clamp}
    \end{equation}
    where $n$ is the mode number within the defect.

  \item \textbf{Defect mode localisation}: The resonant mode is
    localised in a central ``defect'' region of the PnC lattice,
    with exponential decay into the crystal. This minimises radiation
    loss and provides a well-defined effective mass.
\end{enumerate}

\subsection{Complete fabrication specification}

\begin{tcolorbox}[colback=green!5!white, colframe=green!50!black,
  title=\textbf{DFD-OPTIMISED Si$_3$N$_4$ MEMS --- FABRICATION SPEC}]

\begin{center}
\renewcommand{\arraystretch}{1.4}
\begin{tabular}{p{4.5cm}p{3cm}p{5cm}}
\toprule
\textbf{Parameter} & \textbf{Value} & \textbf{Rationale/Source} \\
\midrule
\multicolumn{3}{l}{\textbf{Membrane geometry}} \\
\midrule
Substrate & 200~$\mu$m Si & Standard SOI wafer \\
Membrane material & LPCVD Si$_3$N$_4$ & Stoichiometric; Norcada/in-house \\
Membrane thickness $h$ & 20~nm & Maximise $D_{\rm dil} \propto 1/h^2$ \\
Total chip size & 5~mm $\times$ 5~mm & Fits standard DR sample holders \\
Membrane window & 1~mm $\times$ 1~mm & Set by SQUID pickup coil geometry \\
PnC unit cell & Honeycomb, $a = 150$~$\mu$m & Bandgap centred on $f_0$ \\
PnC hole diameter & 120~$\mu$m & Fill factor 0.65 for wide bandgap \\
Number of PnC periods & $3 \times 3$ around defect & Sufficient for soft clamping \\
Defect size & 300~$\mu$m $\times$ 300~$\mu$m & $\sim 2a$ defect \\
\midrule
\multicolumn{3}{l}{\textbf{Stress and mechanical properties}} \\
\midrule
Tensile stress $\sigma$ & 1.2--1.4~GPa & LPCVD stoichiometric Si$_3$N$_4$ \\
Young's modulus $E$ & 250~GPa & Literature value \\
Density $\rho$ & 3100~kg/m$^3$ & Literature value \\
Poisson's ratio $\nu$ & 0.23 & Literature value \\
\midrule
\multicolumn{3}{l}{\textbf{Resonant mode properties}} \\
\midrule
Mode type & Fundamental defect mode & Localised in central defect \\
Resonant frequency $f_0$ & 1.0--1.5~MHz & Set by defect size and $\sigma$ \\
Effective mass $m_{\rm eff}$ & 0.4--0.8~ng & $\approx 0.4 \times \rho h A_{\rm defect}$ \\
Dilution factor $D_{\rm dil}$ & $\sim 5\ee{5}$ & $\sigma L^2/(E h^2)$ \\
Soft-clamping factor & $\sim 20\times$ & PnC bandgap isolation \\
Total effective dilution & $\sim 10^7$ & $D_{\rm dil} \times D_{\rm soft}$ \\
\midrule
\multicolumn{3}{l}{\textbf{Quality factor}} \\
\midrule
Intrinsic $Q_{\rm int}$ & $\sim 5000$ & Bulk Si$_3$N$_4$ loss at 20~mK \\
$\Qmech$ (room $T$) & $\sim 10^9$ & Demonstrated (multiple groups) \\
$\Qmech$ (4~K) & $>10^{10}$ (projected) & Reduced phonon--phonon loss \\
$\Qmech$ (20~mK) & $>10^{10}$ (projected) & See Section~\ref{sec:q_status} \\
\midrule
\multicolumn{3}{l}{\textbf{Metallisation for SQUID readout}} \\
\midrule
Metal layer & Al (superconducting) & $T_c = 1.2$~K; SQUID-compatible \\
Metal thickness & 15--20~nm & $< 5\%$ mass addition \\
Deposition & e-beam evaporation & Low-stress, uniform \\
Patterning & Liftoff or shadow mask & Avoid plasma damage to Si$_3$N$_4$ \\
\midrule
\multicolumn{3}{l}{\textbf{SQUID coupling}} \\
\midrule
Transduction scheme & Flux modulation & Al-coated membrane as moving \\
 & & element of SQUID pickup loop \\
Gap to pickup coil & 10--50~$\mu$m & Maximise $dB/dx$ \\
Displacement sensitivity & $\sim 10^{-17}$~m/$\sqrt{\rm Hz}$ & Typical dc-SQUID readout \\
\bottomrule
\end{tabular}
\end{center}
\end{tcolorbox}

\subsection{Resonant frequency calculation}

For a square membrane defect of side $a_d$ under tensile stress $\sigma$,
the fundamental mode frequency is:
\begin{equation}
  f_0 = \frac{1}{2a_d}\sqrt{\frac{\sigma}{\rho}}
  = \frac{1}{2 \times 3\ee{-4}}\sqrt{\frac{1.2\ee{9}}{3100}}
  = \frac{1}{6\ee{-4}}\sqrt{3.87\ee{5}}
  = \frac{622}{6\ee{-4}}
  \approx 1.04~\text{MHz}.
  \label{eq:f0_calc}
\end{equation}

This is in the target range. The PnC bandgap must be designed to
encompass $f_0$, typically spanning $0.5f_0$ to $2f_0$ for a
honeycomb lattice with fill factor $\sim 0.65$.

\subsection{Effective mass calculation}

For a soft-clamped defect mode, the effective mass is:
\begin{equation}
  m_{\rm eff} = \eta\,\rho\,h\,a_d^2,
  \label{eq:meff}
\end{equation}
where $\eta \approx 0.4$ is the mode shape factor (the fundamental mode
has slightly lower effective mass than the total defect mass due to the
sinusoidal profile). Thus:
\begin{equation}
  m_{\rm eff} = 0.4 \times 3100 \times 20\ee{-9} \times (3\ee{-4})^2
  = 0.4 \times 3100 \times 20\ee{-9} \times 9\ee{-8}
  = 2.2\ee{-12}~\text{kg}
  \approx 2.2~\text{pg}.
  \label{eq:meff_calc_pg}
\end{equation}

\caution{This is 2.2~pg, not the nanogram regime.} To reach $m_{\rm eff}
\sim 1$~ng, we must scale up the defect size to
$a_d \approx 700$~$\mu$m (with proportionally larger PnC lattice), or
increase the membrane thickness to $h \approx 100$~nm. We present both
options:

\begin{center}
\renewcommand{\arraystretch}{1.3}
\begin{tabular}{lllll}
\toprule
Design & Defect $a_d$ & Thickness $h$ & $m_{\rm eff}$ & $f_0$ \\
\midrule
A: Thin membrane, large defect & 700~$\mu$m & 20~nm & 1.2~ng & 440~kHz \\
B: Thick membrane, std defect & 300~$\mu$m & 100~nm & 1.1~ng & 1.04~MHz \\
C: Intermediate & 500~$\mu$m & 50~nm & 1.2~ng & 620~kHz \\
\midrule
D: Ultra-thin picogram & 300~$\mu$m & 20~nm & 2.2~pg & 1.04~MHz \\
\bottomrule
\end{tabular}
\end{center}

\begin{remark}
Design A maximises the dissipation dilution ($D_{\rm dil} \propto 1/h^2$)
but requires a larger PnC lattice (total membrane $>3$~mm), which is
more fragile and harder to integrate with SQUID pickup geometry.
Design B is more compact and robust but sacrifices $25\times$ in dilution
factor ($(100/20)^2 = 25$). Design C is the recommended compromise for
the Phase~II prototype.
\end{remark}

\subsection{Quality factor: demonstrated vs projected}
\label{sec:q_status}

\begin{theorem}[Si$_3$N$_4$ $\Qmech$ Status: Demonstrated vs Projected]
\label{thm:q_status}
The state of the art for Si$_3$N$_4$ phononic crystal membrane $\Qmech$:

\begin{center}
\renewcommand{\arraystretch}{1.3}
\begin{tabular}{p{3.5cm}rrrl}
\toprule
System & $\Qmech$ & Mass & Temperature & Status \\
\midrule
Soft-clamped PnC, & $1.1\ee{9}$ & $\sim$10~ng & 300~K &
  \textbf{Demonstrated} \\
\quad Tsaturyan et al.\ (2017) & & & & \\
PnC membrane, & $8\ee{8}$ & $\sim$4~ng & 300~K &
  \textbf{Demonstrated} \\
\quad Ghadimi et al.\ (2018) & & & & \\
Trampoline resonator, & $3.6\ee{8}$ & $\sim$50~ng & 300~K &
  \textbf{Demonstrated} \\
\quad Reinhardt et al.\ (2016) & & & & \\
String resonator, & $8\ee{8}$ & $\sim$1~pg & 300~K &
  \textbf{Demonstrated} \\
\quad Ghadimi et al.\ (2018) & & & & \\
PnC membrane at 10~mK, & $>10^{10}$ & $\sim$10~ng & 10~mK &
  \textbf{Projected} \\
\quad (extrapolated) & & & & \\
\midrule
\multicolumn{5}{l}{\textit{Metallised (SQUID-compatible) variants:}} \\
\midrule
Al-coated Si$_3$N$_4$, & $\sim 10^8$ & $\sim$10~ng & 300~K &
  \textbf{Demonstrated} \\
\quad (various groups) & & & & \\
Al-coated at $<1$~K, & $>10^{9}$ & $\sim$ng & $<1$~K &
  \textbf{Projected} \\
\quad (expected, not yet published) & & & & \\
\bottomrule
\end{tabular}
\end{center}

\textbf{Critical assessment}:
\begin{enumerate}[nosep]
  \item $\Qmech \geq 10^9$ at nanogram mass and room temperature:
    \textbf{DEMONSTRATED} in bare (unmetallised) Si$_3$N$_4$ membranes.
  \item $\Qmech \geq 10^9$ at nanogram mass and millikelvin temperature:
    \textbf{NOT YET DEMONSTRATED}, but projected with high confidence.
    The dominant loss mechanisms at room temperature (phonon--phonon
    scattering, thermoelastic damping) freeze out at cryogenic
    temperatures, so $\Qmech$ is expected to \emph{increase} upon
    cooling. The remaining loss (TLS surface loss) typically improves
    by $10\times$--$100\times$ at millikelvin temperatures for thin
    dielectric films.
  \item $\Qmech \geq 10^9$ for \textbf{metallised} Si$_3$N$_4$:
    \textbf{NOT YET DEMONSTRATED at any temperature}. The Al
    metallisation (15--20~nm) introduces additional loss from the
    metal film (surface oxide TLS, electron--phonon coupling at $T > T_c$).
    At $T < T_c^{\rm Al} = 1.2$~K, the Al becomes superconducting and
    electron--phonon loss vanishes, but TLS loss in the native Al$_2$O$_3$
    persists.

    \honest{The metallisation $Q$ penalty is the primary risk for the
    DFD MEMS programme. If Al metallisation degrades $\Qmech$ to
    $\sim 10^8$ (the room-temperature demonstrated value), the Phase~II
    sensitivity degrades by $10\times$ from $1.5\ee{-15}$ to
    $\sim 1.5\ee{-14}$---still a factor of $100\times$ below the
    SQMS bound and therefore still valuable.}
\end{enumerate}
\end{theorem}

\subsection{Fabrication process flow}

\begin{enumerate}
  \item \textbf{Wafer selection}: 200~$\mu$m Si(100) wafer, double-side
    polished, high-resistivity ($>10$~k$\Omega\cdot$cm).

  \item \textbf{Si$_3$N$_4$ deposition}: LPCVD at 780--850$^\circ$C,
    stoichiometric SiH$_2$Cl$_2$/NH$_3$ ratio, targeting $h = 20$--100~nm
    with stress $\sigma = 1.2$~GPa ($\pm 5\%$). Thickness uniformity
    $< 2\%$ across wafer.

  \item \textbf{PnC patterning}: E-beam lithography (for prototype) or
    DUV stepper (for production). PMMA resist, $\sim 200$~nm.

  \item \textbf{PnC etch}: Reactive ion etching (RIE) with CHF$_3$/O$_2$
    chemistry. Etch through Si$_3$N$_4$ only; stop on Si.

  \item \textbf{Membrane release}: KOH backside etch (70$^\circ$C,
    30\% w/w) through Si wafer, stopping on Si$_3$N$_4$ front layer.
    Alternatively, XeF$_2$ vapour etch for smaller windows.

  \item \textbf{Critical point drying}: CO$_2$ supercritical drying to
    prevent membrane collapse during release.

  \item \textbf{Al metallisation} (if SQUID readout): e-beam evaporation
    of 15--20~nm Al through shadow mask (avoids resist processing on
    released membrane). Alternatively, deposit Al before release and
    protect during KOH etch.

  \item \textbf{Characterisation}: Optical interferometric $\Qmech$
    measurement at room temperature to verify $>10^8$. Cryogenic
    characterisation in dilution refrigerator.
\end{enumerate}

\subsection{P2+P11 sensitivity recalculation with specification}
\label{sec:sensitivity_recalc}

Using the Design C parameters ($m_{\rm eff} = 1.2$~ng,
$\Qmech = 10^9$, $f_0 = 620$~kHz) in the W3i5 hybrid formula:

\begin{equation}
  \lam_{\rm hybrid} = \sqrt{
  \frac{2\muG c^4\,m_{\rm eff}\,\Opsi\,\hbar}
  {G_{00}^2\,U_{\rm src}^2\,m^2\,\Qmech^2\,d^{-2}\,T_{\rm int}}
  }.
\end{equation}

With $m_{\rm eff} = 1.2\ee{-9}$~kg (slightly larger than W3i6
assumption of $10^{-9}$~kg) and $\Qmech = 10^9$:

\begin{equation}
  \lam_{\rm min}^{\rm Design\,C} \approx \sqrt{1.2}\times 1.5\ee{-15}
  = 1.6\ee{-15}.
\end{equation}

If the metallisation penalty limits $\Qmech$ to $10^8$:
\begin{equation}
  \lam_{\rm min}^{\rm degraded} \approx 1.6\ee{-14}.
\end{equation}

\begin{finding}[Si$_3$N$_4$ MEMS Fabrication Specification: Validated]
\label{find:si3n4_spec}
The complete fabrication specification is technically feasible using
established MEMS processing. The critical unknowns are:
\begin{enumerate}[nosep]
  \item Metallisation $Q$ penalty at millikelvin temperatures
    (requires measurement);
  \item PnC bandgap integrity after Al deposition (requires FEM
    simulation and measurement);
  \item SQUID coupling efficiency for 620~kHz mode (requires
    transducer design).
\end{enumerate}
The worst-case (metallisation-limited) sensitivity of
$\lam < 1.6\ee{-14}$ is still $100\times$ below the SQMS bound,
confirming the P2+P11 MEMS path as robust.
\end{finding}


%=============================================================================
\section{Nb Cavity Optimisation for Phase~I SQMS 4-Cavity Array}
\label{sec:nb_optimisation}
%=============================================================================

\subsection{Context: Phase~I requirements}

The Phase~I SQMS 4-cavity array (W3i6 P8/P11, Configuration 1) requires
four TESLA 9-cell cavities operated at Fermilab SQMS with:
\begin{itemize}[nosep]
  \item Maximum $Q_0$ at the operating temperature;
  \item Minimum $Q_0$ fluctuation ($\sigma_Q/Q_0$) over integration
    periods of $10^5$--$10^6$~s;
  \item Minimum systematic Q-factor drift (temperature, magnetic field,
    RF level dependence).
\end{itemize}

The passive bound scales as $\lam \propto \sqrt{\sigma_Q/Q_0}$, so
both $Q_0$ and its stability matter equally.

\subsection{Optimal Nb surface preparation: the N-doping recipe}

\begin{tcolorbox}[colback=green!5!white, colframe=green!50!black,
  title=\textbf{RECOMMENDED Nb PREPARATION PROTOCOL FOR PHASE~I}]

\textbf{Step 1: Bulk Nb selection}
\begin{itemize}[nosep]
  \item RRR $\geq$ 300 (residual resistivity ratio); higher RRR
    provides lower residual resistance at low $T$.
  \item Single large-grain or fine-grain from the same ingot
    (for cavity-to-cavity reproducibility in the 4-cavity array).
  \item Tantalum content $< 500$~ppm (Ta substitutional defects
    act as pair-breaking centres).
  \item Hydrogen degassing: 800$^\circ$C vacuum anneal for 3~h
    to remove dissolved hydrogen (prevents Q-disease from NbH$_x$
    precipitation).
\end{itemize}

\textbf{Step 2: Surface preparation}
\begin{itemize}[nosep]
  \item \textbf{Electropolishing (EP)}: Remove 150~$\mu$m of damaged
    surface layer. Standard EP solution: HF(48\%):H$_2$SO$_4$(96\%)
    = 1:9 by volume, 30$^\circ$C, 20~V. This establishes the
    atomically smooth surface ($R_a < 0.5$~$\mu$m) needed for
    minimum surface resistance.
  \item \textbf{Nitrogen doping}: The critical step for maximising $Q_0$.
\end{itemize}

\textbf{Step 3: N-doping (medium temperature)}

The ``2/6'' or ``2/0/2'' Fermilab recipe:
\begin{enumerate}[nosep]
  \item Evacuate furnace to $<10^{-6}$~Torr.
  \item Ramp to 800$^\circ$C at 5$^\circ$C/min in vacuum.
  \item Hold at 800$^\circ$C for 20~min in vacuum (annealing phase).
  \item Inject 25~mTorr N$_2$ for 2~min (``2/6'' $\equiv$ 2~min N$_2$
    + 6~min anneal; ``2/0/2'' $\equiv$ 2~min N$_2$ + 0~min + 2~min
    N$_2$).
  \item Continue 800$^\circ$C anneal for 6~min under vacuum
    (diffusion phase).
  \item Cool at $\leq 10^\circ$C/min to room temperature.
\end{enumerate}

\textbf{Result}: N atoms diffuse $\sim 1$--2~$\mu$m into the Nb
surface, creating a nitrogen-enriched layer that:
\begin{itemize}[nosep]
  \item Reduces the density of two-level system (TLS) defects in the
    native Nb$_2$O$_5$ oxide;
  \item Introduces a controlled mean-free-path reduction that shifts
    the BCS surface resistance minimum to lower temperatures;
  \item Achieves the ``anti-Q-slope'': $Q_0$ \emph{increases} with
    accelerating gradient in the 5--15~MV/m range, which is the
    operating regime for passive measurements.
\end{itemize}

\textbf{Step 4: Final EP}
\begin{itemize}[nosep]
  \item Remove 5~$\mu$m by light EP (removes surface nitrides formed
    during N-doping while preserving the bulk nitrogen diffusion profile).
  \item Rinse with ultrapure water (18.2~M$\Omega\cdot$cm).
\end{itemize}

\textbf{Step 5: Low-temperature bake (120$^\circ$C bake)}
\begin{itemize}[nosep]
  \item 120$^\circ$C for 48~h in vacuum ($<10^{-7}$~Torr).
  \item \textbf{Purpose}: Reduces the high-field Q-slope (less critical
    for passive measurements at low field) but also modifies the
    Nb$_2$O$_5$/NbO$_x$ oxide bilayer, reducing TLS loss by
    transforming amorphous Nb$_2$O$_5$ into partially crystalline NbO$_x$.
  \item The 120$^\circ$C bake is the single most important step for
    \emph{stability} of $Q_0$ over long timescales.
\end{itemize}

\textbf{Step 6: Assembly and cooldown}
\begin{itemize}[nosep]
  \item Clean-room assembly ($<$~Class~100).
  \item Slow cooldown through $T_c$ at $<0.5$~K/min to maximise
    flux expulsion (Meissner effect). Fast cooldown traps ambient
    flux as Abrikosov vortices.
  \item External magnetic shielding to $<1$~nT
    (mu-metal + superconducting Pb sheet).
  \item Operating temperature: 2.0~K (superfluid He bath) for Phase~I
    baseline; 20~mK (dilution refrigerator) for Phase~I upgrade.
\end{itemize}
\end{tcolorbox}

\subsection{Expected $Q_0$ performance}

\begin{center}
\renewcommand{\arraystretch}{1.3}
\begin{tabular}{lrrr}
\toprule
Configuration & $T_{\rm op}$ & $Q_0$ (expected) & $\sigma_Q/Q_0$ \\
\midrule
N-doped, 2~K, superfluid He & 2.0~K & $3\ee{10}$--$5\ee{10}$ &
  $\sim 10^{-4}$ (microphonics-limited) \\
N-doped, 2~K, optimised shielding & 2.0~K & $5\ee{10}$ &
  $\sim 10^{-5}$ (TLS + flux) \\
N-doped, 20~mK, DR & 20~mK & $>2\ee{11}$ &
  $\sim 10^{-6}$ (TLS-limited) \\
N-doped, 20~mK, 300$^\circ$C baked & 20~mK & $>5\ee{11}$ &
  $\sim 10^{-7}$ (projected) \\
\bottomrule
\end{tabular}
\end{center}

\begin{remark}
The 300$^\circ$C bake (``mid-$T$ bake'') is an emerging protocol that
eliminates the Nb$_2$O$_5$ amorphous oxide entirely, replacing it with
a thin ($\sim 2$~nm) NbO$_x$ suboxide layer. This has been shown to
reduce TLS noise by $>10\times$ in superconducting qubit resonators.
For the DFD Phase~I upgrade, implementing a 300$^\circ$C bake after
N-doping could push $Q_0 > 5\ee{11}$ at 20~mK, tightening the passive
bound by $\sqrt{2.5}\approx 1.6\times$.
\end{remark}

\subsection{The $Q_0$ anomaly signature: indirect evidence of $\psi$-coupling}
\label{sec:q0_anomaly}

\begin{theorem}[$Q_0$ Anomaly as Indirect $\psi$-Evidence]
\label{thm:q0_anomaly}
If the DFD $\psi$-field couples to the EM field inside the cavity
via Process~A, the coupling introduces an additional dissipation channel
with a characteristic spectral signature. Specifically:

\begin{enumerate}
  \item \textbf{The $\psi$-induced surface resistance}: Process~A
    transfers energy from the EM mode to the $\psi$-mode at rate:
    \begin{equation}
      \Gamma_{\rm A} = (\lambda-1)^2\,\omega_{\rm EM}^2\,
      |\psiEM_{\rm zpf}|^2\,\Heff^2\,\frac{V_{\rm cav}}{c^2},
    \end{equation}
    contributing an additional surface resistance:
    \begin{equation}
      R_{\psi} = \frac{\Gamma_{\rm A}\,\mu_0\omega}{Q_{\rm geom}}
      = (\lambda-1)^2\,\frac{\mu_0\omega^3}{Q_{\rm geom}}\,
      \frac{\hbar}{2\Meff\Opsi}\,\Heff^2\,V_{\rm cav}/c^2.
      \label{eq:R_psi}
    \end{equation}

  \item \textbf{Temperature dependence}: Unlike BCS surface resistance
    ($R_{\rm BCS} \propto e^{-\Delta/(k_BT)}\,\omega^2/T$), the
    $\psi$-induced resistance is temperature-independent (it couples
    to the zero-point $\psi$-mode, not to thermal excitations):
    \begin{equation}
      R_\psi(T) = R_\psi^0 = \text{const.}
    \end{equation}

  \item \textbf{The anomaly signature}: At sufficiently low temperature
    ($T \ll \Tc$), $R_{\rm BCS} \to 0$ exponentially, leaving only the
    residual resistance $R_{\rm res} = R_{\rm impurity} + R_{\rm TLS}
    + R_\psi$. If the standard residual terms ($R_{\rm impurity}$,
    $R_{\rm TLS}$) can be characterised and subtracted:
    \begin{equation}
      \boxed{
      Q_0^{-1}(T\to 0) - Q_{\rm res}^{-1} = \frac{R_\psi}{\mu_0\omega G}
      = (\lambda-1)^2 \cdot K_{\rm cav},
      }
      \label{eq:anomaly}
    \end{equation}
    where $K_{\rm cav}$ is a calculable cavity-dependent constant.

  \item \textbf{The specific measurement}: Measure $Q_0$ as a function
    of temperature from $\Tc$ down to 20~mK. Fit the BCS + residual
    model:
    \begin{equation}
      Q_0^{-1}(T) = \frac{A}{T}\omega^2 n_s(T) e^{-\Delta/(k_BT)}
      + R_{\rm res}/(\mu_0\omega G).
    \end{equation}
    If, after subtracting the BCS contribution and the
    known residual ($R_{\rm impurity}$ from RRR, $R_{\rm TLS}$ from
    power dependence), a \textbf{frequency-dependent residual} remains
    with:
    \begin{equation}
      \delta(Q_0^{-1}) \propto \omega^2 \quad\text{(not $\omega^{-1}$ or const.)},
      \label{eq:freq_dependence}
    \end{equation}
    this would be consistent with Process~A and inconsistent with
    standard residual resistance mechanisms (which are
    frequency-independent or scale as $\omega^{-1}$).
\end{enumerate}

\honest{This anomaly would be \emph{suggestive} but not \emph{proof}
of $\psi$-coupling. An unexplained $\omega^2$ residual resistance could
also arise from surface pair-breaking, paramagnetic impurities, or
unknown TLS physics. It would constitute \textbf{indirect, circumstantial
evidence} that motivates the Phase~II MEMS experiment, which provides a
fundamentally different detection signature.}

\textbf{The specific measurement for Phase~I}: If the 4-cavity array
achieves $\sigma_{Q_0}/Q_0 < 10^{-6}$ at 20~mK and 1.3~GHz, measure
each cavity also at 650~MHz (half-wave mode) and 2.6~GHz (third
harmonic). Compare $R_{\rm res}(\omega)$. If $R_{\rm res} \propto
\omega^2$ with a coefficient consistent with the DFD prediction for
$\lam \sim 10^{-9}$, this is the anomaly.
\end{theorem}

\begin{finding}[Phase~I Multi-Frequency Strategy]
\label{find:multi_freq}
The Phase~I 4-cavity array should measure $Q_0$ at multiple harmonics
(0.65, 1.3, 2.6~GHz) of each cavity. The TESLA 9-cell geometry supports
these modes. A residual resistance with $R_{\rm res} \propto \omega^2$
after BCS subtraction would constitute the first indirect evidence of
psi-EM coupling, with significance depending on the ability to exclude
mundane explanations.
\end{finding}


%=============================================================================
\section{Novel Material Candidates: Assessment}
\label{sec:novel}
%=============================================================================

\subsection{Scope}

W3i6 considered Nb, Nb$_3$Sn, MgB$_2$, YBCO, Si$_3$N$_4$, Al, diamond,
and CNT/PhC nanobeams. The W3i7 mandate asks whether four additional
material classes offer exceptional passive FOM:
\begin{enumerate}[nosep]
  \item Superconducting nanowires;
  \item Topological surface states;
  \item Superfluid $^3$He;
  \item Dilute Bose gas (BEC).
\end{enumerate}

\subsection{Candidate 1: Superconducting nanowires}
\label{sec:nanowires}

\textbf{The physics}: Superconducting nanowire single-photon detectors
(SNSPDs) are NbN, NbTiN, WSi, or MoSi nanowires ($w \sim 50$--100~nm,
$h \sim 5$--10~nm) that detect single photons via the formation of a
resistive hotspot. Could the extreme confinement of supercurrents in a
nanowire geometry enhance the psi-EM overlap integral $\Heff$?

\textbf{Analysis}:
\begin{itemize}[nosep]
  \item The overlap integral $\Heff$ depends on the EM field distribution
    in the superconductor. In a nanowire ($w < \lambdaL$), the
    current is approximately uniform across the cross-section, but the
    total volume of superconducting material is tiny
    ($V_{\rm wire} \sim w \times h \times L \sim 10^{-18}$~m$^3$ per wire).
  \item For passive detection, we need high $Q_0$. Nanowire resonators
    have been demonstrated as kinetic inductance detectors (KIDs) with
    $Q_{\rm int} \sim 10^6$--$10^7$, which is $10^4\times$ lower than
    bulk Nb cavities.
  \item The passive FOM $= Q_0 \cdot \Heff$. Even if $\Heff$ is
    enhanced $100\times$ by current confinement, the $10^4\times$
    lower $Q_0$ makes nanowires net $100\times$ worse than Nb cavities.
\end{itemize}

\begin{finding}[Superconducting Nanowires: Not Competitive]
\label{find:nanowires}
\killed{Superconducting nanowires are not competitive for passive
DFD detection.} The $Q$ penalty from nanoscale geometry (surface loss,
kinetic inductance dissipation) overwhelms any enhancement in the
psi-EM overlap. The passive FOM is $\sim 10^{16}$, compared with
$10^{20}$ for bulk Nb.

\textbf{Possible niche}: Nanowires could serve as single-photon
detectors in an \emph{active} scheme where psi-to-photon conversion
events are counted individually. However, the active detection gap
($10^{12}$--$10^{17}$) makes this approach nonviable in the near term.
\end{finding}

\subsection{Candidate 2: Topological surface states}
\label{sec:topological}

\textbf{The physics}: Topological insulators (Bi$_2$Se$_3$,
Bi$_2$Te$_3$, Sb$_2$Te$_3$) and topological semimetals (Cd$_3$As$_2$,
Na$_3$Bi, TaAs) host protected surface states with linear (Dirac)
dispersion. When proximitised with a superconductor, these surface
states may host Majorana fermions. Could the topological protection
mechanism yield an enhanced or anomalous psi-coupling?

\textbf{Analysis}:
\begin{itemize}[nosep]
  \item The DFD psi-field couples to EM stress-energy via the
    $(\lambda-1)$ vertex. This coupling is determined by the EM field
    amplitude and the material's dielectric/superconducting response,
    not by the topological properties of the band structure.
  \item Topological protection prevents backscattering of surface
    currents, which is relevant for transport but does not affect the
    stored EM energy or Q-factor in a resonator geometry.
  \item Topological insulator thin films on superconducting substrates
    achieve $Q \sim 10^4$--$10^5$ in microwave resonators, limited by
    the non-superconducting bulk and surface disorder.
  \item The passive FOM for a TI-based resonator:
    $\FOMpassive \sim 10^5 \times \Heff \sim 10^{10}$--$10^{12}$,
    far below Nb.
\end{itemize}

\begin{finding}[Topological Surface States: Not Competitive]
\label{find:topological}
\killed{Topological surface states do not provide an enhanced
psi-coupling mechanism.} The topological protection operates in
momentum space (preventing backscattering) and does not affect the
real-space EM field distribution that determines the psi-EM overlap.
The Q-factor limitation of TI-based resonators ($10^4$--$10^5$)
eliminates them from passive detection consideration.

\textbf{Theoretical note}: If DFD were to predict an anomalous coupling
between the $\psi$-field and the Berry phase of topological surface
states (a higher-order effect beyond Process~A), this would require
a modification of the DFD field equation. No such modification has
been proposed.
\end{finding}

\subsection{Candidate 3: Superfluid $^3$He}
\label{sec:he3}

\textbf{The physics}: Superfluid $^3$He (below $T_c \approx 2.5$~mK at
saturated vapour pressure, or $\sim 1$~mK at zero pressure) is a
p-wave superfluid with extraordinarily low intrinsic dissipation.
Torsional oscillators and vibrating wire resonators in superfluid $^3$He
have demonstrated acoustic quality factors $Q_{\rm acoustic} > 10^{10}$
at temperatures below 100~$\mu$K.

\textbf{How it could detect psi}: Superfluid $^3$He supports
well-defined acoustic modes (zero sound, first sound) with exceptionally
low damping. If the $\psi$-field couples to the mass-energy density of
the $^3$He (which it must, through $\psigrav$), then psi-field
oscillations would drive acoustic modes of the superfluid. A
torsional oscillator or vibrating wire immersed in superfluid $^3$He
acts as a high-$Q$ detector of density fluctuations.

\textbf{The passive FOM for superfluid $^3$He}:
\begin{itemize}[nosep]
  \item $Q_{\rm acoustic}$: $>10^{10}$ (demonstrated in torsional
    oscillators at Lancaster, Helsinki, Cornell).
  \item The psi-coupling mechanism: the $\psi$-field perturbs the local
    gravitational potential, which shifts the density of the $^3$He:
    $\delta\rho/\rho = -\psi/c^2$. This density shift couples to
    acoustic modes with an effective coupling constant:
    \begin{equation}
      g_{\rm He3} = \frac{\rho_{\rm He3}\,V_{\rm mode}}{c^2}
      \,\omega_{\rm acoustic}\,\psi_{\rm zpf}.
    \end{equation}
  \item \textbf{Critical difference from SRF}: This is a
    \emph{gravitational} coupling ($\psigrav$), not an EM coupling
    ($\psiEM$). The $(\lambda-1)$ factor does not enter. This means
    superfluid $^3$He is sensitive to the gravitational psi-field, which
    has amplitude $\psigrav \approx 0.047$, not the EM component
    ($\psiEM \approx 9.6\ee{-6}$).
\end{itemize}

\begin{theorem}[Superfluid $^3$He: New Passive Detection Candidate]
\label{thm:he3}
Superfluid $^3$He torsional oscillators offer a qualitatively different
passive detection pathway that couples to $\psigrav$ rather than
$\psiEM$:

\begin{enumerate}
  \item \textbf{Coupling mechanism}: Gravitational---the psi-field
    perturbs the local density, driving acoustic modes.
  \item \textbf{Advantage}: The gravitational coupling is $\sim 5000\times$
    stronger than the EM coupling ($\psigrav/\psiEM \approx 0.047/9.6\ee{-6}
    \approx 4900$).
  \item \textbf{Q-factor}: $>10^{10}$, comparable to Si$_3$N$_4$ MEMS
    and exceeding SRF Nb by $\sim 50\times$.
  \item \textbf{Mode mass}: A superfluid $^3$He cell of volume
    $V = 1$~cm$^3$ has effective acoustic mode mass
    $m_{\rm eff} = \rho V/2 = 81$~kg/m$^3 \times 10^{-6}$~m$^3$ / 2
    $= 4\ee{-5}$~kg $= 40$~$\mu$g.
    This is much heavier than the Si$_3$N$_4$ MEMS, which penalises
    sensitivity.
  \item \textbf{Net passive sensitivity}: The sensitivity scales as
    $\sqrt{m_{\rm eff}}/Q$. Compared to Si$_3$N$_4$ ($m = 1$~ng,
    $Q = 10^9$):
    \begin{equation}
      \frac{\lam_{\rm He3}}{\lam_{\rm Si3N4}}
      = \frac{\sqrt{4\ee{-5}}/10^{10}}{\sqrt{10^{-9}}/10^{9}}
      = \frac{6.3\ee{-3}/10^{10}}{3.2\ee{-5}/10^{9}}
      = \frac{6.3\ee{-13}}{3.2\ee{-14}} \approx 20.
    \end{equation}
    Superfluid $^3$He is $\sim 20\times$ less sensitive than Si$_3$N$_4$
    MEMS for the EM coupling channel.
\end{enumerate}

\honest{However, the superfluid $^3$He couples to $\psigrav$, not
$\psiEM$. If the goal is to constrain $|\lambda-1|$ (the EM-psi coupling
constant), superfluid $^3$He provides no direct constraint on $\lambda$.
It constrains the gravitational $\psi$-field properties, which are
already constrained by solar system tests and MOND phenomenology.}

\textbf{Verdict}: Superfluid $^3$He is an \textbf{interesting new
candidate for gravitational psi-detection} but \textbf{does not
constrain $|\lambda-1|$} and therefore does not advance the primary DFD
detection programme (which seeks to measure $\lambda$). It belongs in
Track~3 (monitoring) as a complementary probe.

\caution{If DFD is correct and $\psigrav$ has the predicted properties,
a superfluid $^3$He acoustic oscillator at $T < 100$~$\mu$K could
detect psi-gravitational effects that are independent of $\lambda$.
This would be confirmatory evidence for DFD even without constraining
the EM coupling.}
\end{theorem}

\subsection{Candidate 4: Dilute Bose gas (BEC)}
\label{sec:bec}

\textbf{The physics}: A dilute Bose-Einstein condensate
($^{87}$Rb, $^{23}$Na, $^{7}$Li) at nanokelvin temperatures is a
macroscopic quantum state with extremely low density
($\rho \sim 10^{14}$~atoms/cm$^3$, or $\rho \sim 10^{-11}$~kg/m$^3$).
BECs support collective modes (phonons, Bogoliubov excitations) with
Q-factors limited by three-body loss and thermal depletion.

\textbf{Analysis for DFD passive detection}:
\begin{itemize}[nosep]
  \item \textbf{Coupling}: Like superfluid $^3$He, a BEC couples to
    $\psigrav$ gravitationally (density perturbation), not to $\psiEM$.
  \item \textbf{Q-factor}: Collective mode lifetimes in BECs are
    typically 0.1--10~s, giving $Q \sim \omega\tau \sim 10^3$--$10^5$
    at typical mode frequencies of 100~Hz--10~kHz. This is far below
    SRF cavities or Si$_3$N$_4$.
  \item \textbf{Mode mass}: A typical BEC of $10^6$ atoms has
    $m = 10^6 \times 1.4\ee{-25}$~kg $= 1.4\ee{-19}$~kg $= 0.14$~fg.
    This is extremely low mass, which would be advantageous, but the
    low $Q$ overwhelms the mass advantage.
  \item \textbf{Passive FOM}: $\FOMMEMS^{\rm BEC} \propto Q^{3/2}/m
    \sim (10^4)^{3/2}/(10^{-19}) = 10^{25}$. This looks impressive,
    but the FOM must be evaluated in terms of actual displacement
    sensitivity. The BEC displacement noise is dominated by quantum
    shot noise and atom number fluctuations, not thermal noise, and
    the readout sensitivity for BEC collective modes is orders of
    magnitude worse than SQUID displacement detection.
\end{itemize}

\begin{finding}[Dilute Bose Gas: Not Competitive]
\label{find:bec}
\killed{Dilute Bose gases are not competitive for DFD passive detection.}
The low Q-factor ($\sim 10^4$) and poor readout sensitivity relative to
SQUID-coupled solid-state resonators eliminate BECs from consideration.
Furthermore, BECs couple to $\psigrav$ (not $\psiEM$), so they cannot
constrain $|\lambda-1|$.

\textbf{Atom interferometry note}: While the BEC collective modes are
not useful, \emph{atom interferometers} using cold atoms (but not
necessarily BECs) are the leading technology for gravitational
psi-detection in the P8 communications receiver. This is a different
application (single-particle interferometry, not collective mode
detection).
\end{finding}

\subsection{Novel materials summary}

\begin{center}
\renewcommand{\arraystretch}{1.3}
\begin{tabular}{p{3.5cm}p{2cm}p{2cm}p{2cm}p{3cm}}
\toprule
Material & Couples to & $Q$ & $\lam$ constraint? & Verdict \\
\midrule
SC nanowires & $\psiEM$ & $10^6$--$10^7$ & Yes (weak) &
  \killed{Not competitive} \\
Topological SS & $\psiEM$ & $10^4$--$10^5$ & Yes (very weak) &
  \killed{Not competitive} \\
Superfluid $^3$He & $\psigrav$ & $>10^{10}$ & \textbf{No} ($\psigrav$ only) &
  \caution{Track 3 monitor} \\
Dilute Bose gas & $\psigrav$ & $\sim 10^4$ & \textbf{No} ($\psigrav$ only) &
  \killed{Not competitive} \\
\bottomrule
\end{tabular}
\end{center}

\begin{finding}[Novel Material Assessment: No Track~1/2 Changes]
\label{find:novel_summary}
None of the four novel materials displaces Nb (passive) or Si$_3$N$_4$
(MEMS) from the primary programme. Superfluid $^3$He is added to
Track~3 as a complementary gravitational psi-probe.
\end{finding}


%=============================================================================
\section{Cryogenic Engineering: Cost and Power Budget}
\label{sec:cryo}
%=============================================================================

\subsection{Phase~I: SQMS 4-cavity array}

\subsubsection{Option A: 2~K operation (superfluid He-4)}

The SQMS facility at Fermilab operates TESLA cavities at 2~K using
superfluid helium. The cryogenic parameters:

\begin{center}
\renewcommand{\arraystretch}{1.3}
\begin{tabular}{lrl}
\toprule
Parameter & Value & Notes \\
\midrule
Number of cavities & 4 & TESLA 9-cell \\
Static heat load per cavity & $\sim$2~W at 2~K & Radiation, conduction, RF \\
Total 2~K heat load & $\sim$10~W & Including transfer lines \\
Carnot efficiency at 2~K & $1/150$ & $T_{\rm cold}/(T_{\rm hot}-T_{\rm cold})$ \\
Plant efficiency factor & $\sim$0.25 & Practical vs Carnot \\
Wall-plug power & $10/(1/150\times 0.25) = 6$~kW & Cryoplant compressor \\
LHe consumption & $\sim$20~L/day & If non-recirculating \\
\midrule
Capital cost (if existing) & \$0 & Fermilab SQMS has the plant \\
Capital cost (if new) & \$2--3M & Complete 2~K cryoplant \\
Operating cost & $\sim$\$50k/yr & Electricity + LHe makeup \\
\bottomrule
\end{tabular}
\end{center}

\subsubsection{Option B: 20~mK operation (dilution refrigerator)}

For the Phase~I upgrade to 20~mK, a dilution refrigerator (DR) must
cool the cavity assembly:

\begin{center}
\renewcommand{\arraystretch}{1.3}
\begin{tabular}{lrl}
\toprule
Parameter & Value & Notes \\
\midrule
DR cooling power at 20~mK & $\sim$200~$\mu$W & Large-capacity DR \\
Heat load at 20~mK & $<100$~$\mu$W & 4 cavities + wiring \\
DR model & Bluefors XLD1000 class & Or Oxford Instruments Proteox \\
Pre-cooling stages & 50~K, 4~K, still, CP, MXC & Standard dry DR \\
4~K stage power & $\sim$1.5~W (pulse tube) & Per cooler \\
Number of pulse tubes & 2 & For 4-cavity mass \\
Total wall-plug power & $\sim$50~kW & Compressors + electronics \\
\midrule
Capital cost (DR) & \$1.5--2.5M & Large custom DR \\
Capital cost (custom) & \$3--4M & With cavity integration \\
Operating cost & $\sim$\$100k/yr & Electricity + maintenance \\
\bottomrule
\end{tabular}
\end{center}

\begin{remark}
A single TESLA 9-cell cavity has mass $\sim$25~kg and length $\sim$1~m.
Four cavities ($\sim$100~kg) require a custom DR with an unusually large
sample space. The Bluefors XLD series can accommodate large payloads
but may need modification for 4-cavity loading. An alternative is to
operate one cavity at a time in a standard DR and combine results
(losing simultaneous cross-correlation capability).
\end{remark}

\subsection{Phase~II: P2+P11 hybrid with $\mu$g MEMS}

\begin{center}
\renewcommand{\arraystretch}{1.3}
\begin{tabular}{lrl}
\toprule
Parameter & Value & Notes \\
\midrule
Cryogenic requirement & 2~K (source cavity) + & Source SRF at 2~K; \\
 & 20~mK (MEMS detector) & MEMS in DR \\
Source cavity cooling & Existing 2~K He plant & Fermilab or equivalent \\
MEMS DR & Standard DR & Single-cavity or standalone \\
 & (Bluefors LD400 class) & MEMS mass $\sim$~mg \\
DR cooling power & 400~$\mu$W at 20~mK & Standard capacity \\
Wall-plug power & $\sim$30~kW (DR) + & \\
 & 6~kW (He plant) $= 36$~kW & \\
\midrule
Capital cost (DR) & \$0.5--1M & Standard DR \\
Capital cost (He plant) & \$0 (existing) & \\
Operating cost & $\sim$\$80k/yr & \\
\bottomrule
\end{tabular}
\end{center}

\subsection{Phase~III: 256-cavity coherent array}

\begin{tcolorbox}[colback=red!5!white, colframe=red!50!black,
  title=\textbf{PHASE~III CRYOGENIC COST: FACILITY-CLASS INVESTMENT}]

\begin{center}
\renewcommand{\arraystretch}{1.3}
\begin{tabular}{lrl}
\toprule
Parameter & Value & Notes \\
\midrule
Number of cavities & 256 & TESLA 9-cell, N-doped Nb \\
Static heat load & $\sim$500~W at 2~K & 256 $\times$ 2~W \\
Carnot + efficiency & 500~W $\times$ 600 & = 300~kW wall-plug \\
Dynamic RF load & $\sim$100~W at 2~K & Low-field passive operation \\
Total 2~K load & $\sim$600~W & \\
Wall-plug power & $\sim$360~kW & $\approx 0.4$~MW \\
\midrule
If operated at 20~mK: & & \\
\quad Pre-cooling & 256 cavities $\times$ 25~kg & = 6400~kg payload \\
\quad DR capacity & Multiple large DRs & 10--20 XLD-class units \\
\quad DR wall-plug & $\sim$50~kW $\times$ 10 & = 500~kW \\
\quad Total wall-plug & $\sim$0.5--1~MW & At 20~mK \\
\midrule
Alternative: 4.2~K (Nb$_3$Sn) & $\sim$600~W at 4.2~K & More efficient \\
\quad Wall-plug & $\sim$50~kW & $\sim$8$\times$ cheaper than 2~K \\
\midrule
\textbf{Capital costs:} & & \\
Cryoplant (2~K, new) & \$15--25M & Large He refrigerator \\
DRs (20~mK, multiple) & \$15--30M & Custom multi-DR system \\
Cryogenic distribution & \$10--20M & Transfer lines, valves, controls \\
\textbf{Total cryo capital} & \textbf{\$40--75M} & \\
\midrule
\textbf{Operating cost} & \$1--3M/yr & Electricity + LHe + maintenance \\
\bottomrule
\end{tabular}
\end{center}

\textbf{Comparison}: The Phase~III cryogenic budget (\$40--75M capital,
$\sim$0.5--1~MW operating) is comparable to:
\begin{itemize}[nosep]
  \item The European XFEL cryoplant (2~K, 5~kW, \$50M);
  \item The LCLS-II cryoplant at SLAC (2~K, 2~kW, \$40M);
  \item The ESS Lund cryoplant (2~K, 8~kW, \$60M).
\end{itemize}
Phase~III is unambiguously a facility-class project.
\end{tcolorbox}

\subsection{Impact on programme economics}
\label{sec:economics}

\begin{finding}[Cryogenic Cost Impact on Programme Economics]
\label{find:cryo_economics}
\begin{enumerate}
  \item \textbf{Phase~I (2~K, 4 cavities)}: Cryogenic costs are
    negligible---\$0 capital (Fermilab existing), $\sim$\$50k/yr
    operating. The Phase~I budget of \$2--4M is dominated by personnel,
    not cryogenics.

  \item \textbf{Phase~I upgrade (20~mK, 4 cavities)}: DR cost of
    \$1.5--4M is significant---comparable to the rest of the Phase~I
    budget. This upgrade should be a decision gate: only if Phase~I at
    2~K identifies a systematic floor above $10^{-10}$ that requires
    lower temperature to circumvent.

  \item \textbf{Phase~II (MEMS hybrid)}: Cryogenic costs are modest
    ($\sim$\$0.5--1M for a standard DR). The Phase~II budget of
    \$3--6M is dominated by personnel and MEMS development.

  \item \textbf{Phase~III (256-cavity)}: \textbf{Cryogenics dominates
    the budget}. The \$40--75M cryo cost is 50--60\% of the total
    Phase~III estimate of \$80--120M. This makes Phase~III a
    facility-class decision requiring institutional commitment from a
    national laboratory.

  \item \textbf{Alternative Phase~III (ng MEMS)}: The P2+P11 hybrid
    with nanogram Si$_3$N$_4$ MEMS achieves comparable sensitivity
    ($\lam \sim 10^{-15}$) with a \$5--10M budget and $\sim$\$1M
    cryogenic cost. \textbf{This is the cost-optimal path to deep
    reach}, confirming the W3i6 finding that Track~2 (MEMS) is a
    competitive alternative to Track~1 Phase~3 (256-cavity array).
\end{enumerate}
\end{finding}

\subsection{Power budget summary}

\begin{center}
\renewcommand{\arraystretch}{1.3}
\begin{tabular}{lrrr}
\toprule
Phase & Wall-plug power & Capital cryo & Operating cryo/yr \\
\midrule
I (2~K, existing) & 6~kW & \$0 & \$50k \\
I (20~mK upgrade) & 50~kW & \$1.5--4M & \$100k \\
II (MEMS hybrid) & 36~kW & \$0.5--1M & \$80k \\
III (256 $\times$ 2~K) & 400~kW & \$15--25M & \$1--2M \\
III (256 $\times$ 20~mK) & 1000~kW & \$40--75M & \$2--3M \\
III alt (ng MEMS) & 36~kW & \$0.5--1M & \$80k \\
\bottomrule
\end{tabular}
\end{center}


%=============================================================================
\section{Updated Material Property Table}
\label{sec:property_table}
%=============================================================================

Incorporating W3i7 additions (superfluid $^3$He, refined Si$_3$N$_4$
parameters):

\begin{table}[H]
\centering
\caption{Comprehensive material properties for DFD detection, updated W3i7.
Properties at $T_{\rm op} = 20$~mK unless noted.}
\label{tab:properties_w3i7}
\renewcommand{\arraystretch}{1.2}
\begin{tabular}{p{2.2cm}rrrrrr}
\toprule
Property & Nb & Nb$_3$Sn & Si$_3$N$_4$ & Al & $^3$He & BEC \\
\midrule
Role & SRF passive & SRF alt. & MEMS & MEMS alt. & Novel & Novel \\
Couples to & $\psiEM$ & $\psiEM$ & $\psiEM$ & $\psiEM$ & $\psigrav$ & $\psigrav$ \\
$T_c$ (K) & 9.25 & 18.3 & --- & 1.2 & 0.0025 & --- \\
$Q$ (best) & $2\ee{11}$ & $2\ee{10}$ & $10^{10}$ & $10^7$ & $>10^{10}$ & $10^4$ \\
Min mass & 25~kg & 25~kg & 1~pg--1~ng & 1~ng & 40~$\mu$g & 0.1~fg \\
$\FOMpassive$ & $\mathbf{10^{20}}$ & $10^{19}$ & --- & $10^{19}$ & $10^{15}$* & $10^{10}$ \\
$\FOMMEMS$ & $10^{7}$ & --- & $\mathbf{10^{18}}$ & $10^{10}$ & --- & $10^{12}$** \\
\bottomrule
\end{tabular}

\smallskip
\small
*$^3$He FOM not directly comparable: couples to $\psigrav$, not $\psiEM$.\\
**BEC FOM inflated by ultra-low mass; readout sensitivity limits practical use.
\end{table}


%=============================================================================
\section{Ranked Findings}
\label{sec:ranked}
%=============================================================================

\begin{enumerate}
  \item \textbf{[RANK 1 --- NEW] Si$_3$N$_4$ MEMS Fabrication Spec
    Validated.}
    The complete fabrication specification for a soft-clamped PnC
    Si$_3$N$_4$ membrane is technically feasible. Design C (500~$\mu$m
    defect, 50~nm thick, $f_0 = 620$~kHz, $m_{\rm eff} = 1.2$~ng)
    is the recommended Phase~II prototype. $\Qmech \geq 10^9$ at
    nanogram mass and room temperature is \textbf{demonstrated};
    at millikelvin temperature it is \textbf{projected with high
    confidence}. The metallisation $Q$ penalty is the primary risk;
    even the worst case ($\Qmech = 10^8$) gives $\lam < 1.6\ee{-14}$,
    which is $100\times$ below the SQMS bound.
    (Theorem~\ref{thm:q_status}, Finding~\ref{find:si3n4_spec})

  \item \textbf{[RANK 2 --- NEW] Optimal Nb Preparation Protocol for
    Phase~I.}
    Medium-$T$ N-doping (800$^\circ$C, 2/0/2 recipe), 5~$\mu$m EP,
    120$^\circ$C $\times$ 48~h bake, sub-nT shielding, slow cooldown
    through $T_c$. Expected $Q_0 > 2\ee{11}$ at 20~mK with
    $\sigma_Q/Q_0 \sim 10^{-6}$. Optional 300$^\circ$C bake upgrade
    could reach $Q_0 > 5\ee{11}$.
    (Section~\ref{sec:nb_optimisation})

  \item \textbf{[RANK 3 --- NEW] $Q_0$ Anomaly Signature Defined.}
    A residual resistance scaling as $R_{\rm res} \propto \omega^2$
    (measured across cavity harmonics at 0.65, 1.3, 2.6~GHz) after
    BCS subtraction would constitute circumstantial evidence of
    psi-EM coupling. This is the Phase~I multi-frequency measurement
    strategy. Evidence is indirect; MEMS experiment needed for
    confirmation.
    (Theorem~\ref{thm:q0_anomaly}, Finding~\ref{find:multi_freq})

  \item \textbf{[RANK 4 --- NEW] Superfluid $^3$He: Gravitational
    Psi-Probe.}
    Superfluid $^3$He ($Q > 10^{10}$) couples to $\psigrav$ with
    $\sim 5000\times$ stronger amplitude than $\psiEM$, but does
    \textbf{not} constrain $|\lambda-1|$. Added to Track~3 as
    complementary gravitational psi-probe. Not a substitute for
    the primary Nb/Si$_3$N$_4$ programme.
    (Theorem~\ref{thm:he3})

  \item \textbf{[RANK 5 --- NEW] Superconducting Nanowires and
    Topological Surface States: Eliminated.}
    Neither provides competitive passive FOM. Nanowires:
    $Q \sim 10^6$--$10^7$, $\FOMpassive \sim 10^{16}$ (vs $10^{20}$
    for Nb). Topological: $Q \sim 10^4$--$10^5$, no enhanced psi-coupling.
    (Findings~\ref{find:nanowires}, \ref{find:topological})

  \item \textbf{[RANK 6 --- NEW] Dilute Bose Gas: Eliminated.}
    $Q \sim 10^4$, couples to $\psigrav$ only, poor readout sensitivity.
    (Finding~\ref{find:bec})

  \item \textbf{[RANK 7 --- NEW] Cryogenic Budget Validates Phase~I/II
    Economics.}
    Phase~I cryo cost is negligible (\$0 capital at Fermilab).
    Phase~II cryo cost is modest (\$0.5--1M). Phase~III cryogenics
    (\$40--75M) dominates the 256-cavity budget and makes it
    facility-class. The ng-MEMS alternative (Phase~III alt) achieves
    comparable reach at $1/80$ the cryo cost.
    (Finding~\ref{find:cryo_economics})

  \item \textbf{[RANK 8 --- CONFIRMED] Three-Track Roadmap Reinforced.}
    The W3i6 three-track structure (Passive/Nb, MEMS/Si$_3$N$_4$,
    Exotic/monitoring) is strengthened by: (a) validated fabrication
    spec for Track~2; (b) validated Nb recipe for Track~1; (c) cryo
    cost analysis confirming ng-MEMS as the cost-optimal deep-reach
    path.
\end{enumerate}


%=============================================================================
\section{Cross-Reference Updates to Other Agents}
\label{sec:cross_ref}
%=============================================================================

\begin{itemize}[nosep]
  \item \textbf{To P8/P11 Agent}: The Phase~I multi-frequency strategy
    (Sec.~\ref{sec:q0_anomaly}) should be incorporated into the Phase~I
    experimental protocol. The $R_{\rm res} \propto \omega^2$ signature
    is a falsifiable prediction that can be tested with the 4-cavity
    array at no additional cost.

  \item \textbf{To P2 Agent}: The Si$_3$N$_4$ MEMS specification
    (Sec.~\ref{sec:si3n4_spec}) provides concrete parameters for the
    P2+P11 hybrid sensitivity calculation. The metallisation $Q$ penalty
    introduces a factor-of-10 uncertainty in Phase~II sensitivity
    ($1.6\ee{-15}$ vs $1.6\ee{-14}$).

  \item \textbf{To Cosmo Agent}: Superfluid $^3$He (Sec.~\ref{sec:he3})
    couples to $\psigrav$, not $\psiEM$. If the Cosmo agent has
    predictions for $\psigrav$ oscillation modes or spectral features
    at acoustic frequencies ($\sim$~kHz), these should be cross-referenced
    with the $^3$He torsional oscillator sensitivity.

  \item \textbf{To P14/P15 Agent}: The superconducting nanowire
    assessment (Sec.~\ref{sec:nanowires}) confirms that nanowire-based
    single-photon detection of psi-to-photon conversion is not viable
    given the active detection gap. No change to the P15 (active
    detection killed) assessment.
\end{itemize}


%=============================================================================
\section{Open Items for W3i8}
\label{sec:open}
%=============================================================================

\begin{enumerate}[nosep]
  \item \textbf{Si$_3$N$_4$ metallisation $Q$ measurement}: The single
    most important experimental measurement for Track~2 is $\Qmech$ of
    a 15--20~nm Al-coated soft-clamped PnC Si$_3$N$_4$ membrane at
    $T < 1$~K. Request from experimental groups (Delft, Copenhagen,
    NIST) or propose as a Phase~II precursor experiment.

  \item \textbf{SQUID transducer design}: Detailed design of the
    flux-to-displacement transducer for the Al-coated Si$_3$N$_4$
    membrane. Required deliverable: displacement sensitivity
    ($\text{m}/\sqrt{\text{Hz}}$) vs frequency for the SQUID readout
    geometry at 20~mK.

  \item \textbf{Phase~I multi-frequency protocol}: Develop the
    BCS-subtraction analysis pipeline for extracting $R_{\rm res}(\omega)$
    from multi-harmonic $Q_0$ measurements. Quantify the expected
    $R_\psi$ for $\lam = \lambdabound$ and determine the measurement
    precision needed to detect or exclude it.

  \item \textbf{Superfluid $^3$He scoping}: Estimate the sensitivity
    of a Lancaster/Helsinki-type torsional oscillator to $\psigrav$
    perturbations at acoustic frequencies. Determine whether this
    provides competitive constraints on gravitational sector DFD
    parameters.

  \item \textbf{Nb cavity aging}: What is the long-term ($>1$~year)
    $Q_0$ stability of N-doped Nb cavities? If $Q_0$ degrades over
    time (oxide regrowth, hydrogen re-dissolution), this limits the
    integration time for passive bounds. Investigate whether periodic
    re-baking is needed.

  \item \textbf{Phase~III facility siting}: If Phase~III proceeds,
    which national laboratory offers the best existing cryogenic
    infrastructure? Candidates: Fermilab (SQMS), DESY (XFEL), ESS
    (Lund). Cost differential of co-locating vs building new.
\end{enumerate}


\end{document}
